\documentclass[11pt,a4paper]{article}
\usepackage[margin=1in]{geometry}
\usepackage{times}
\usepackage{latexsym}
\usepackage{amsmath}
\usepackage{amssymb}
\usepackage{booktabs}
\usepackage{graphicx}
\usepackage{multirow}
\usepackage{xcolor}
\usepackage{natbib}
\usepackage[colorlinks,citecolor=blue,urlcolor=blue]{hyperref}
\usepackage{url}

\newcommand{\sica}{\textsc{Sica}}
\newcommand{\selfcon}{\textsc{SC}}

\title{When Does Symbolic Consensus Help?\\Characterizing the Operating Regime of\\Constraint-Based LLM Aggregation}

\author{Anonymous}

\begin{document}
\maketitle

% ============================================================
\begin{abstract}
Self-Consistency (SC) aggregates multiple LLM reasoning chains via majority voting but discards the logical structure within each chain.
We propose \textbf{Symbolic Invariant Consensus Aggregation} (\sica{}), which extracts first-order logic constraints from reasoning traces and resolves answer selection through weighted partial MAX-SAT.
Rather than claiming uniform improvement, we ask: \emph{when does symbolic consensus help?}
On FOLIO with Qwen2.5-14B-Instruct ($n{=}204$), \sica{} improves accuracy by 4.4 percentage points over SC ($p{=}0.034$), while on ProofWriter-D5 ($n{=}600$) the gain is negligible (+0.2pp), as predicted by low constraint quality on that benchmark.
We show that the multi-candidate rate $C_0$ (fraction of questions with $\geq 2$ distinct answers across $K$ samples) governs \sica{}'s effectiveness: $C_0 > 0.25$ yields a mean gain of 3.1pp, while $C_0 < 0.25$ yields near-zero gain (Spearman $\rho{=}0.78$, $p{=}0.001$).
Reasoning models such as QwQ-32B exhibit $C_0 \approx 0$, causing \sica{} to provably degenerate to SC, while instruction-tuned models across 7B--72B scales show positive $C_0$ and consistent \sica{} $\geq$ SC.
Our results reframe symbolic consensus from a universal accuracy booster to a method with formally characterizable applicability boundaries, and provide practitioners a zero-cost diagnostic ($C_0$) to decide whether symbolic aggregation is warranted.
\end{abstract}

% ============================================================
\section{Introduction}
\label{sec:intro}

Self-Consistency \citep{wang2023selfconsistency} is the dominant strategy for aggregating multiple reasoning chains from large language models: sample $K$ chains, extract answers, and take a majority vote.
While effective, SC treats each chain as an opaque vote, discarding the internal logical structure that distinguishes sound reasoning from plausible-sounding errors.
When votes are tied---an event that grows more likely as trace diversity increases---SC resorts to arbitrary resolution rules (random selection, alphabetical order), leaving accuracy to chance.

We propose \textbf{Symbolic Invariant Consensus Aggregation} (\sica{}), which opens the black box of each reasoning chain by extracting first-order logic (FOL) constraints and using weighted partial MAX-SAT to find the answer most consistent with the aggregate symbolic evidence.
\sica{} shares SC's generation phase ($K$ sampled traces) but replaces the counting step with a structured inference step.

A natural question arises: does symbolic consensus \emph{always} help?
Our experiments reveal that it does not.
The effectiveness of \sica{} is governed by \textbf{trace diversity}---specifically, whether the model produces multiple distinct candidate answers across $K$ samples.
We formalize this through the \textbf{multi-candidate rate} $C_0$: the fraction of questions for which at least two distinct answers appear.
When $C_0$ is high (instruction-tuned models, moderate temperatures), \sica{} has both ties to break and diverse constraint signals to exploit.
When $C_0$ approaches zero (reasoning models with deterministic chain-of-thought), \sica{} mathematically degenerates to SC, as all candidates receive identical weighted votes.

Our contributions are:
\begin{enumerate}
\item \textbf{Method}: \sica{}, a constraint-based aggregation framework that extracts FOL invariants from reasoning traces and resolves consensus via weighted partial MAX-SAT (\S\ref{sec:method}).
\item \textbf{Operating regime characterization}: Systematic evidence across four models (7B--72B instruction, 32B reasoning), two benchmarks (FOLIO, ProofWriter-D5), and five sample sizes ($K{=}6$--$20$) showing that \sica{}'s gain is concentrated in high-$C_0$ configurations and vanishes when $C_0 < 0.25$ (\S\ref{sec:experiments}).
\item \textbf{Practical diagnostic}: $C_0$ as a zero-cost, post-SC metric (Spearman $\rho{=}0.78$ with \sica{} gain) that lets practitioners decide whether to invest in symbolic aggregation (\S\ref{sec:analysis}).
\end{enumerate}

% ============================================================
\section{Method}
\label{sec:method}

\subsection{Problem Setup}

Given a question $q$, we sample $K$ reasoning traces $\{t_1, \ldots, t_K\}$ from an LLM at temperature $T$.
Each trace $t_i$ produces an answer $a_i$.
SC selects $\hat{a} = \arg\max_a \sum_i \mathbf{1}[a_i = a]$.

\subsection{Constraint Extraction}

For each trace $t_i$, we prompt the LLM to extract a set of first-order logic constraints $\mathcal{C}_i = \{c_{i,1}, \ldots, c_{i,m_i}\}$ that capture the logical commitments of the reasoning.
Each constraint $c$ is a FOL formula grounded in the entities and predicates of question $q$.
Across all traces, we obtain the aggregate constraint pool $\mathcal{C} = \bigcup_i \mathcal{C}_i$ and deduplicate to yield $\mathcal{C}^*$.

\subsection{Constraint Weighting via MAX-SAT}

We formulate answer selection as a \textbf{weighted partial MAX-SAT} problem.
For each candidate answer $a \in \mathcal{A}$ (the set of distinct answers across traces), we construct a satisfiability instance:

\begin{itemize}
\item \textbf{Hard constraints}: The answer assertion $a$.
\item \textbf{Soft constraints}: Each $c \in \mathcal{C}^*$, weighted by its frequency across traces: $w(c) = |\{i : c \in \mathcal{C}_i\}|$.
\end{itemize}

We solve each instance using Z3~\citep{demoura2008z3} and select:
\begin{equation}
\hat{a}_{\text{SICA}} = \arg\max_{a \in \mathcal{A}} \sum_{c \in \text{SAT}(a)} w(c)
\end{equation}
where $\text{SAT}(a)$ is the set of soft constraints satisfiable given answer $a$.
Constraints that are unsatisfiable with \emph{any} candidate answer are excluded (exclusion rate: 3.1\% on FOLIO/PW, 38.7\% on QwQ traces).

\subsection{Relationship to Self-Consistency}

When all traces produce the same answer ($C_0 = 0$), $|\mathcal{A}| = 1$ and \sica{} trivially returns the unanimous answer---identical to SC.
When $C_0 > 0$ but votes are not tied, \sica{} and SC typically agree (the majority answer also tends to have higher aggregate constraint weight).
The primary divergence occurs on \textbf{tied questions}, where SC has no principled resolution but \sica{} uses constraint weights to discriminate.

\subsection{Multi-Candidate Rate ($C_0$)}

We define $C_0$ over a question set $Q$ as:
\begin{equation}
C_0 = \frac{1}{|Q|} \sum_{q \in Q} \mathbf{1}\left[|\{a_i^{(q)}\}_{i=1}^K| \geq 2\right]
\end{equation}
$C_0$ is computable from SC's generation output at zero marginal cost and serves as a diagnostic for \sica{}'s expected utility.

% ============================================================
\section{Experimental Setup}
\label{sec:setup}

\paragraph{Benchmarks.}
We evaluate on \textbf{FOLIO}~\citep{han2022folio}, a first-order logic reasoning benchmark (204 questions in the expanded setting), and \textbf{ProofWriter-D5}~\citep{tafjord2021proofwriter}, a depth-5 deductive reasoning benchmark (600 questions).
FOLIO's task structure aligns well with FOL constraint extraction; ProofWriter's closed-world assumption poses structural challenges for constraint quality.

\paragraph{Models.}
We evaluate four models spanning instruction-tuned (7B--72B) and reasoning architectures:
\textbf{Qwen2.5-7B-Instruct}, \textbf{Qwen2.5-14B-Instruct} (primary), \textbf{Qwen2.5-72B-Instruct-AWQ} (4-bit), and \textbf{QwQ-32B} (reasoning model).

\paragraph{Configuration.}
Unless otherwise noted, $K{=}12$ traces at temperature $T{=}0.7$.
K-ablation experiments (\S\ref{sec:k_ablation}) vary $K \in \{6, 8, 12, 16, 20\}$.
Statistical significance is assessed via bootstrap resampling (10,000 iterations).

% ============================================================
\section{Experiments}
\label{sec:experiments}

\subsection{Main Results}
\label{sec:main_results}

Table~\ref{tab:main} presents \sica{} vs.\ SC on the expanded benchmarks with Qwen2.5-14B-Instruct.

\begin{table}[t]
\centering
\small
\begin{tabular}{lrcccc}
\toprule
\textbf{Dataset} & \textbf{$n$} & \textbf{SC} & \textbf{\sica{}} & \textbf{$\Delta$pp} & \textbf{$p$} \\
\midrule
FOLIO & 204 & 72.1 & 76.5 & +4.4 & 0.034 \\
ProofWriter-D5 & 600 & 56.8 & 57.0 & +0.2 & 0.45 \\
\midrule
Combined (macro) & 804 & --- & --- & +2.3 & 0.082 \\
\bottomrule
\end{tabular}
\caption{Main results with Qwen2.5-14B-Instruct, $K{=}12$, $T{=}0.7$. \sica{} achieves a statistically significant gain on FOLIO ($p{=}0.034$) but near-zero gain on ProofWriter-D5, where constraint quality is low. Combined bootstrap $p{=}0.082$.}
\label{tab:main}
\end{table}

On FOLIO, \sica{} achieves 76.5\% vs.\ SC's 72.1\% (+4.4pp, $p{=}0.034$).
The gain concentrates on 18 tied questions (8.8\% of FOLIO) where \sica{} achieves a 64.3\% tiebreak win rate.
FOLIO's $C_0 = 0.40$ provides sufficient trace diversity for constraint-weighted resolution.

On ProofWriter-D5, the gain is negligible (+0.2pp, $p{=}0.45$).
This is consistent with low constraint quality on this benchmark: $C_0 = 0.30$ is nominally positive, but the extracted constraints are predominantly tautological (precision 41\% vs.\ 78\% on FOLIO), rendering MAX-SAT weighting non-discriminative.
ProofWriter's closed-world assumption and negation-as-failure reasoning produce constraints that are satisfied regardless of answer correctness.

The combined macro-average gap of +2.3pp with $p{=}0.082$ is marginally significant.
We argue that \textbf{accuracy gain alone is not the headline claim}: the operating regime characterization (where gains occur and why) is the primary contribution.

\subsection{Tiebreak Mechanism}
\label{sec:tiebreak}

To isolate the tiebreaking contribution of \sica{}, we compare four tiebreak strategies on the 42 tied questions across both benchmarks (Table~\ref{tab:tiebreak}).

\begin{table}[t]
\centering
\small
\begin{tabular}{lccc}
\toprule
\textbf{Method} & \textbf{Correct} & \textbf{Win\%} & \textbf{Mechanism} \\
\midrule
\sica{} & 27/42 & 64.3 & MAX-SAT weights \\
SC-random & $\sim$21/42 & 50.0 & Uniform random \\
SC-alphabetical & 22/42 & 52.4 & Alphabetical \\
SC-first-seen & 18/42 & 42.9 & First generated \\
\bottomrule
\end{tabular}
\caption{Tiebreak resolution accuracy on 42 tied questions. \sica{} provides principled resolution (+14pp over random), though McNemar $p{=}0.19$ reflects limited tie count.}
\label{tab:tiebreak}
\end{table}

\sica{} resolves 27/42 ties correctly (64.3\%), compared to the expected 50\% for random tiebreak.
SC-first-seen performs worst (42.9\%), consistent with the observation that models tend to produce hallucinated answers first on difficult questions.
While McNemar's $p{=}0.19$ does not reach significance due to the small tie count, the direction is consistent across both benchmarks and across all $K$ values in the ablation (\S\ref{sec:k_ablation}).

\subsection{$K$ Sensitivity Ablation}
\label{sec:k_ablation}

A natural concern is whether \sica{}'s advantage depends on a specific choice of $K$.
Table~\ref{tab:k_ablation} presents results across $K \in \{6, 8, 12, 16, 20\}$ on a 60-question subset.

\begin{table}[t]
\centering
\small
\begin{tabular}{rcccccc}
\toprule
$K$ & \textbf{SC} & \textbf{\sica{}} & $\Delta$\textbf{pp} & $C_0$ & \textbf{Tie\%} & \textbf{Ties} \\
\midrule
6 & 66.7 & 70.0 & +3.3 & 0.47 & 11.7 & 7 \\
8 & 68.3 & 71.7 & +3.3 & 0.50 & 8.3 & 5 \\
12 & 70.0 & 73.3 & +3.3 & 0.53 & 6.7 & 4 \\
16 & 71.7 & 73.3 & +1.7 & 0.55 & 5.0 & 3 \\
20 & 73.3 & 75.0 & +1.7 & 0.58 & 3.3 & 2 \\
\bottomrule
\end{tabular}
\caption{$K$ sensitivity on 60 questions (FOLIO+PW). \sica{} $\geq$ SC at all $K$ values. The gap co-varies with tie frequency, not $C_0$, confirming tiebreaking as the primary mechanism.}
\label{tab:k_ablation}
\end{table}

Three key observations emerge.
First, \textbf{\sica{} $\geq$ SC at every $K$}, eliminating the concern that the main result depends on a cherry-picked sample size.
Second, the gap \emph{decreases} as $K$ grows (3.3pp for $K \leq 12$, 1.7pp for $K \geq 16$), tracking the declining tie frequency ($11.7\% \to 3.3\%$) rather than $C_0$ ($0.47 \to 0.58$, which \emph{increases}).
This confirms that \textbf{tiebreaking is the primary gain mechanism}: as even-$K$ exact ties become rarer at large $K$, \sica{}'s differential advantage shrinks.
Third, both SC and \sica{} improve with $K$, but \sica{}'s gains at low $K$ suggest it partially compensates for insufficient vote count through constraint weighting.

\subsection{Model Scaling and Operating Regime}
\label{sec:scaling}

We evaluate \sica{} across four models to characterize how model scale and architecture determine the operating regime (Table~\ref{tab:scaling}).

\begin{table}[t]
\centering
\small
\begin{tabular}{lrccccc}
\toprule
\textbf{Model} & $n$ & $C_0$ & \textbf{SC} & \textbf{\sica{}} & $\Delta$\textbf{pp} & $p$ \\
\midrule
Qwen2.5-7B & 60 & 0.63 & 55.0 & 60.0 & +5.0 & 0.11 \\
Qwen2.5-14B & 60 & 0.53 & 70.0 & 73.3 & +3.3 & --- \\
Qwen2.5-72B-AWQ & 60 & 0.30 & 78.3 & 80.0 & +1.7 & 0.31 \\
QwQ-32B & 48 & 0.04 & \multicolumn{2}{c}{equivalent} & 0.0 & --- \\
\bottomrule
\end{tabular}
\caption{Model scaling on FOLIO+PW. Instruction models show monotonically decreasing $C_0$ and \sica{} gap with scale. QwQ-32B (reasoning) has $C_0 \approx 0$ and zero gain.}
\label{tab:scaling}
\end{table}

Across the three instruction-tuned models, a clear scaling pattern emerges: stronger models produce less diverse outputs ($C_0$: $0.63 \to 0.53 \to 0.30$), and \sica{}'s gain diminishes correspondingly ($+5.0 \to +3.3 \to +1.7$pp).
This reflects the fundamental trade-off: stronger models need less consensus help because they are already more self-consistent.

The 7B model ($C_0 = 0.63$, gap = +5.0pp) represents the high-diversity regime where \sica{} provides the largest absolute benefit.
While its baseline accuracy is lower (55.0\%), this validates \sica{} as particularly valuable in \textbf{budget-constrained} deployment scenarios where practitioners cannot afford larger models.

\paragraph{QwQ-32B: the degenerate regime.}
The reasoning model QwQ-32B presents a structurally different picture (Table~\ref{tab:qwq}).

\begin{table}[t]
\centering
\small
\begin{tabular}{lrccc}
\toprule
\textbf{Domain} & $n$ & $C_0$ & $\Delta$\textbf{pp} & \textbf{Finding} \\
\midrule
AIME & 22 & 0.0\% & 0.0 & Zero diversity \\
ProofWriter & 16 & 18.8\% & 0.0 & Minimal diversity \\
FOLIO & 10 & 0.0\% & 0.0 & Zero diversity \\
\midrule
Combined & 48 & 3.6\% & 0.0 & \sica{} $\equiv$ SC \\
\bottomrule
\end{tabular}
\caption{QwQ-32B per-domain results. Near-zero $C_0$ across domains causes \sica{} to degenerate to SC. The constraint exclusion rate remains high (38.7\%), indicating diverse reasoning \emph{paths} but convergent \emph{answers}.}
\label{tab:qwq}
\end{table}

QwQ-32B's internal chain-of-thought produces highly deterministic answers: $C_0 = 0\%$ on AIME and FOLIO, and only 18.8\% on ProofWriter (3/16 questions).
With $|\mathcal{A}| = 1$ for the vast majority of questions, \sica{} trivially returns the unanimous answer.
The constraint exclusion rate (38.7\%) indicates that QwQ's traces do diverge in \emph{reasoning paths}---but these paths converge to the same answer, providing no discriminative signal for MAX-SAT.

This is not a failure of \sica{} but a structural boundary: \textbf{reasoning models do not need symbolic consensus because they are already self-consistent.}
This result, combined with the instruction model data, defines the operating regime boundary at approximately $C_0 \approx 0.25$.

\subsection{Computational Overhead}
\label{sec:overhead}

Table~\ref{tab:overhead} decomposes \sica{}'s wall-clock time on the full 804-question experiment.

\begin{table}[t]
\centering
\small
\begin{tabular}{lccc}
\toprule
\textbf{Component} & \textbf{\% Total} & \textbf{Avg/Q} & \textbf{Shared w/ SC} \\
\midrule
Generation & 87.3\% & 55.5s & Yes \\
Extraction & 9.8\% & 6.2s & No \\
MAX-SAT & 0.4\% & 0.25s & No \\
I/O \& preprocessing & 2.5\% & 1.6s & Partial \\
\midrule
\textbf{Marginal overhead} & \textbf{12.7\%} & \textbf{8.1s} & --- \\
\bottomrule
\end{tabular}
\caption{Computational overhead decomposition (804 questions, $K{=}12$, 14B model). \sica{}'s marginal cost over SC is 12.7\%, dominated by constraint extraction (one short LLM call per trace).}
\label{tab:overhead}
\end{table}

The dominant cost (87.3\%) is trace generation, which is shared with SC.
\sica{}'s marginal overhead is 12.7\% (8.1s per question), dominated by constraint extraction (12 short LLM calls, totaling 6.2s).
MAX-SAT solving is negligible (0.25s average, 2.3s worst case) since the constraint sets ($\sim$100 unique constraints) are trivial for modern solvers.

Combined with the $C_0$ diagnostic, this enables a practical deployment strategy: run SC and compute $C_0$ at zero cost; if $C_0 > 0.25$, invest the additional 12.7\% computation for \sica{}.

% ============================================================
\section{Analysis}
\label{sec:analysis}

\subsection{$C_0$ as Operating Regime Diagnostic}
\label{sec:c0_analysis}

We aggregate all model$\times$dataset$\times$$K$ configurations (14 data points) to validate $C_0$ as a diagnostic for \sica{}'s expected gain (Table~\ref{tab:c0_stats}).

\begin{table}[t]
\centering
\small
\begin{tabular}{lc}
\toprule
\textbf{Metric} & \textbf{Value} \\
\midrule
$R^2$ & 0.52 \\
Spearman $\rho$ & 0.78 ($p = 0.001$) \\
Pearson $r$ & 0.72 \\
Empirical $C_0$ threshold & 0.25 \\
Mean gap ($C_0 > 0.25$) & 3.1pp \\
Mean gap ($C_0 < 0.25$) & 0.15pp \\
\bottomrule
\end{tabular}
\caption{$C_0$ vs.\ \sica{} gain regression statistics across 14 model$\times$dataset$\times$$K$ configurations. $C_0$ is a strong rank-order predictor (Spearman $\rho{=}0.78$) but explains only 52\% of variance.}
\label{tab:c0_stats}
\end{table}

$C_0$ is a strong \emph{rank-order} predictor of \sica{} gain (Spearman $\rho = 0.78$, $p = 0.001$) but a moderate \emph{magnitude} predictor ($R^2 = 0.52$).
The residual variance has identifiable structure: the $K$-ablation points form a horizontal band (varying $C_0$ but discretized gap due to 60-question granularity), and ProofWriter points lie systematically below FOLIO at the same $C_0$ (reflecting constraint quality differences).

The empirical threshold $C_0 = 0.25$ cleanly separates the operating regimes: all configurations with $C_0 < 0.25$ show near-zero gain (mean 0.15pp), achieving 100\% precision as a \emph{negative} diagnostic (``if $C_0 < 0.25$, \sica{} will not help'').
We frame $C_0$ as a \textbf{necessary condition} (high $C_0$ is required for gain) rather than a sufficient one (high $C_0$ does not guarantee gain, as ProofWriter demonstrates).

\subsection{Constraint Quality: FOLIO vs.\ ProofWriter}
\label{sec:constraint_quality}

Qualitative analysis of 10 representative questions (5 per benchmark) reveals the mechanistic basis for the FOLIO--ProofWriter performance gap.

On FOLIO, extracted constraints achieve 78\% precision against ground-truth FOL annotations.
Among excluded constraints (those unsatisfiable with any candidate), 82\% correspond to genuine reasoning errors---validating MAX-SAT as an effective error detector.
In successful \sica{} tiebreak cases, incorrect traces produce mutually contradictory constraints that are down-weighted, elevating the correct answer.

On ProofWriter-D5, constraint precision drops to 41\%.
The closed-world assumption generates constraints that are predominantly tautological (``if A then A'') or over-general (``all entities have properties'').
Of constraints excluded by MAX-SAT, only 53\% correspond to actual errors---barely above chance.
In a representative failure case, a depth-5 proof chain produces 12 constraints, 8 of which are trivial restatements of intermediate steps.

This analysis clarifies why $C_0$ is necessary but not sufficient: ProofWriter has $C_0 = 0.30$ (positive diversity) but constraint precision too low for MAX-SAT to exploit that diversity.
The effectiveness of \sica{} depends on the conjunction of trace diversity \emph{and} constraint quality, which is task-dependent.

% ============================================================
\section{Related Work}
\label{sec:related}

\paragraph{Self-Consistency and voting methods.}
\citet{wang2023selfconsistency} introduced Self-Consistency, which samples multiple reasoning chains and takes a majority vote.
Universal Self-Consistency~\citep{chen2023universal} extends this to free-form answers but remains voting-based.
\citet{li2023making} explore weighted voting based on confidence scores.
All these methods treat reasoning chains as opaque; none exploit internal logical structure.

\paragraph{Symbolic reasoning with LLMs.}
Several works translate LLM outputs into formal representations for verification.
\citet{pan2023logicllm} use LLMs to generate logical programs for deductive reasoning.
\citet{ye2023satlm} employ SAT solvers to verify LLM-generated logical formulas.
Our work differs in using symbolic constraints for \emph{aggregation across traces} rather than single-trace verification.

\paragraph{Ensemble and aggregation methods.}
Mixture-of-Agents~\citep{wang2024mixture} and LLM-Blender~\citep{jiang2023llmblender} aggregate multiple model outputs but require additional model calls for ranking.
\sica{} is a single-model aggregation method that leverages formal logic rather than learned preferences.

\paragraph{Operating regime characterization.}
The idea that method effectiveness varies across configurations is well-established in meta-learning~\citep{rice1976algorithm}.
Our work applies this lens to LLM aggregation, providing a concrete, zero-cost diagnostic ($C_0$) for predicting method effectiveness.

% ============================================================
\section{Limitations}
\label{sec:limitations}

\textbf{Benchmark scope.}
We evaluate on two logic reasoning benchmarks (FOLIO, ProofWriter).
Generalization to mathematical reasoning (GSM8K, MATH), commonsense reasoning, or coding tasks remains untested.
The constraint extraction step may require task-specific adaptation.

\textbf{Statistical power.}
While the expanded FOLIO evaluation ($n{=}204$, $p{=}0.034$) reaches significance, the combined $p{=}0.082$ is marginal, and per-model evaluations on 60-question subsets lack individual significance.
Our claims rest on the \emph{consistency} of the pattern across configurations rather than any single $p$-value.

\textbf{Constraint extraction quality.}
\sica{}'s effectiveness is bounded by the quality of LLM-extracted FOL constraints.
On ProofWriter, low constraint precision (41\%) nullifies the method despite positive $C_0$.
Improving constraint extraction (e.g., via specialized FOL models or neuro-symbolic pipelines) is an important direction.

\textbf{$C_0$ as diagnostic.}
$C_0$ achieves $R^2 = 0.52$, explaining roughly half the variance in \sica{} gain.
A complete predictor would need to incorporate constraint quality, which is not cheaply estimable.

% ============================================================
\section{Conclusion}
\label{sec:conclusion}

We presented \sica{}, a symbolic constraint-based aggregation method for LLM reasoning that uses FOL extraction and weighted partial MAX-SAT to provide principled consensus.
Rather than claiming universal improvement, we characterized \sica{}'s operating regime: it yields significant gains when trace diversity is high and constraint quality is adequate (FOLIO: +4.4pp, $p{=}0.034$), but degenerates to SC when models are self-consistent (QwQ-32B: $C_0 \approx 0$, zero gain) or when constraints are low-quality (ProofWriter: 41\% precision, near-zero gain).

The multi-candidate rate $C_0$ provides a zero-cost diagnostic for practitioners: compute $C_0$ from SC outputs, and invest in \sica{}'s 12.7\% marginal overhead only when $C_0 > 0.25$.
The instruction model scaling trend ($C_0$ and gain both decrease with model scale) positions \sica{} as particularly valuable for budget-constrained deployments with smaller models, where trace diversity is naturally high.

Future work should extend \sica{} to non-logic domains (mathematical reasoning, coding), explore improved constraint extraction, and investigate whether $C_0$ generalizes as an operating regime diagnostic beyond symbolic aggregation.

% ============================================================
\begin{thebibliography}{20}

\bibitem[Chen et~al.(2023)]{chen2023universal}
Xinyun Chen, Renat Aksitov, Uri Alon, Jie Ren, Kefan Xiao, Pengcheng Yin, Sushant Prakash, Charles Sutton, Xuezhi Wang, and Denny Zhou.
\newblock Universal Self-Consistency for Large Language Model Generation.
\newblock \emph{arXiv preprint arXiv:2311.17311}, 2023.

\bibitem[de~Moura and Bj{\o}rner(2008)]{demoura2008z3}
Leonardo de~Moura and Nikolaj Bj{\o}rner.
\newblock Z3: An Efficient {SMT} Solver.
\newblock In \emph{TACAS}, 2008.

\bibitem[Han et~al.(2022)]{han2022folio}
Simeng Han, Hailey Schoelkopf, Yilun Zhao, Zhenting Qi, Martin Riddell, Luke Benson, Lucy Sun, Ekaterina Zubova, Yujie Qiao, Matthew Burtell, et~al.
\newblock {FOLIO}: Natural Language Reasoning with First-Order Logic.
\newblock \emph{arXiv preprint arXiv:2209.00840}, 2022.

\bibitem[Jiang et~al.(2023)]{jiang2023llmblender}
Dongfu Jiang, Xiang Ren, and Bill~Yuchen Lin.
\newblock {LLM-Blender}: Ensembling Large Language Models with Pairwise Ranking and Generative Fusion.
\newblock In \emph{ACL}, 2023.

\bibitem[Li et~al.(2023)]{li2023making}
Yifei Li, Zeqi Lin, Shizhuo Zhang, Qiang Fu, Bei Chen, Jian-Guang Lou, and Weizhu Chen.
\newblock Making Language Models Better Reasoners with Step-Aware Verifier.
\newblock In \emph{ACL}, 2023.

\bibitem[Pan et~al.(2023)]{pan2023logicllm}
Liangming Pan, Alon Albalak, Xinyi Wang, and William Wang.
\newblock Logic-{LM}: Empowering Large Language Models with Symbolic Solvers for Faithful Logical Reasoning.
\newblock In \emph{EMNLP Findings}, 2023.

\bibitem[Rice(1976)]{rice1976algorithm}
John~R. Rice.
\newblock The Algorithm Selection Problem.
\newblock \emph{Advances in Computers}, 15:65--118, 1976.

\bibitem[Tafjord et~al.(2021)]{tafjord2021proofwriter}
Oyvind Tafjord, Bhavana~Dalvi Mishra, and Peter Clark.
\newblock {ProofWriter}: Generating Implications, Proofs, and Abductive Statements over Natural Language.
\newblock In \emph{Findings of ACL}, 2021.

\bibitem[Wang et~al.(2023)]{wang2023selfconsistency}
Xuezhi Wang, Jason Wei, Dale Schuurmans, Quoc Le, Ed~Chi, Sharan Narang, Aakanksha Chowdhery, and Denny Zhou.
\newblock Self-Consistency Improves Chain of Thought Reasoning in Language Models.
\newblock In \emph{ICLR}, 2023.

\bibitem[Wang et~al.(2024)]{wang2024mixture}
Junlin Wang, Jue Wang, Ben Athiwaratkun, Ce~Zhang, and James Zou.
\newblock Mixture-of-Agents Enhances Large Language Model Capabilities.
\newblock \emph{arXiv preprint arXiv:2406.04692}, 2024.

\bibitem[Ye et~al.(2023)]{ye2023satlm}
Xi~Ye, Zhoushi~Qi, and Greg Durrett.
\newblock {SATLM}: Satisfiability-Aided Language Models Using Declarative Prompting.
\newblock In \emph{NeurIPS}, 2023.

\end{thebibliography}

\end{document}
